% Generated draft: requires mathematical and source review.
\documentclass[11pt]{article}
\usepackage[margin=25mm]{geometry}
\usepackage{fontspec}
\setmainfont{Latin Modern Roman}
\usepackage{amsmath,amssymb,mathtools,unicode-math}
\setmathfont{Latin Modern Math}
\usepackage{enumitem,needspace,xurl,hyperref}
\setlength{\parindent}{0pt}
\setlength{\parskip}{5pt}
\setlength{\emergencystretch}{3em}
\newcommand{\N}{\mathbb{N}}
\newcommand{\Z}{\mathbb{Z}}
\newcommand{\Q}{\mathbb{Q}}
\newcommand{\R}{\mathbb{R}}
\newcommand{\C}{\mathbb{C}}
\newcommand{\F}{\mathbb{F}}
\newcommand{\A}{\mathbb{A}}
\newcommand{\PP}{\mathbb{P}}
\newcommand{\E}{\mathbb{E}}
\newcommand{\eps}{\varepsilon}
\newcommand{\dd}{\,\mathrm{d}}
\begin{document}
\section*{UnsolvedMath ID: 30005708}
\textbf{Source problem code:} catalogue entry
\subsection{Definitions and mathematical statement}
A Coxeter system is a group with a presentation
\[
W=\langle S\mid (st)^{m_{st}}=1\text{ whenever }m_{st}<\infty\rangle,
\]
where \(m_{ss}=1\) and \(m_{st}=m_{ts}\in\{2,3,\ldots\}\cup\{\infty\}\) for distinct \(s,t\in S\). Its rank is \(|S|\). The supplied questions do not explicitly require finite rank.

A group \(G\) is residually finite rationally solvable (RFRS) if it admits a descending sequence of finite-index normal subgroups
\[
G=G_0\supseteq G_1\supseteq\cdots,\qquad \bigcap_{j\geq0}G_j=\{1\},
\]
such that
\[
\ker\bigl(G_j\longrightarrow (G_j/[G_j,G_j])\otimes_{\Z}\Q\bigr)\subseteq G_{j+1}
\quad\text{for every }j\geq0.
\]
A group is virtually RFRS if it has a finite-index RFRS subgroup. Its profinite completion is
\[
\widehat G=\varprojlim_{\substack{K\trianglelefteq G\\[G:K]<\infty}}G/K.
\]
Isomorphisms of profinite completions are understood as isomorphisms of topological groups.

One standard convention for the rational Euler characteristic, when a group \(W\) has a torsion-free subgroup \(H\) of finite index with a finite classifying space \(K(H,1)\), is
\[
\chi(W)=\frac{\chi(K(H,1))}{[W:H]}.
\]
This explains the invariant in the question on that domain; the supplied record does not specify the finiteness hypotheses or a convention covering every Coxeter group it might intend.

A group \(G\) is good in the sense of profinite cohomology if, for every finite discrete \(G\)-module \(M\) and every integer \(n\geq0\), the natural comparison map
\[
H^n_{\mathrm{cts}}(\widehat G,M)\longrightarrow H^n(G,M)
\]
is an isomorphism. Here the action on the finite module extends to \(\widehat G\). We use this standard interpretation of the source's otherwise undefined term ``good.''

The principal questions in the supplied clean statement are the following.
\begin{enumerate}
\item Determine conditions under which a Coxeter group is virtually RFRS. This is Question 8, attributed in the supplied original text to Yuri Santos Rego. The supplied dated curation states that virtual RFRS is known for finite-rank Coxeter groups; this is a supplied assessment, not a result verified here. Consequently, the intended remaining scope of the historical question needs clarification.
\item Do there exist Coxeter groups \(W_1,W_2\), with their rational Euler characteristics defined, such that
\[
\widehat W_1\cong\widehat W_2
\qquad\text{and}\qquad
\chi(W_1)\neq\chi(W_2)?
\]
This is Question 9, attributed to Yuri Santos Rego. No finite-rank or other finiteness restriction is explicit in the supplied statement.
\item Let \(N_q\) be the Neretin group described by the source as acting on the \(q\)-regular rooted tree. Are both assertions
\[
H^k(N_q,B_i(N_q))=0
\qquad\text{and}\qquad
H_c^k(E_{\mathcal O}N_q,\Q)=0
\]
true for every cohomological degree \(k\)? These are the separate parts of Question 10, attributed to Thomas Weigel. The supplied text does not define the allowed parameters \(q\), the tree convention, the coefficient object \(B_i(N_q)\), the family \(\mathcal O\), the space \(E_{\mathcal O}N_q\), or the cohomology theory used in the first expression. In particular, the rendering \(B_i\) might conceal a symbol such as \(\operatorname{Bi}\), rather than an indexed family; no range or quantifier for \(i\) is supplied. The displayed expressions retain the supplied clean notation and are formal targets pending these definitions. They do not constitute a fully specified standalone cohomological problem.
\item Does there exist a Coxeter group that is not good? With the interpretation of goodness above, this asks whether there exist a Coxeter group \(W\), a finite discrete \(W\)-module \(M\), and an integer \(n\geq0\) such that
\[
H^n_{\mathrm{cts}}(\widehat W,M)\longrightarrow H^n(W,M)
\]
is not an isomorphism. This is Question 11, attributed to Thomas Weigel. The intended meaning of ``good'' and any rank restriction require confirmation.
\end{enumerate}

The original source excerpt additionally contains Questions 6 and 7, attributed to Roman Sauer, although the clean statement omits them. To preserve this material without silently merging the targets, their content is recorded separately. A TDLC group is a totally disconnected locally compact group. For a compact-open subgroup \(U\leq G\), normalize a Haar measure \(\mu\) by \(\mu(U)=1\). Question 6 asks whether a version of the Atiyah conjecture holds for TDLC groups, or whether there exist such a group, normalization, and degree \(n\) with
\[
b_n^{(2)}(G,\mu)\notin\Q.
\]
No precise candidate version of the Atiyah conjecture is supplied. A lattice \(\Gamma\leq G\) is a discrete subgroup with finite covolume. Question 7 asks whether there is a compactly generated TDLC group with a lattice of irrational covolume. The normalization for this latter question is not repeated explicitly; the preceding compact-open normalization must be confirmed as its intended convention.

The excerpt's Remark 3 states the proportionality relation
\[
b_n^{(2)}(\Gamma)=\operatorname{covol}(\Gamma)\,b_n^{(2)}(G)
\]
for a locally compact group and lattice, and states that non-compactly generated examples with irrational covolume exist. These are supplied background assertions, not independently verified results here; the measure convention and applicable hypotheses require checking. Remark 4 attributes to Kielak, under its reference label [2], an equivalence between virtual fibering and \(b_1^{(2)}(G)=0\) for finitely generated virtually RFRS groups. The record does not define virtual fibering, identify [2], or supply the full theorem hypotheses. This remark is preserved as a source claim and is not asserted here as a theorem in that generality.
\subsection{Short English statement}
For permanent UnsolvedMath ID \(30005708\), the supplied record asks about virtual RFRS for Coxeter groups, whether profinite completion can fail to determine their rational Euler characteristic, two vanishing questions for Neretin groups, and the existence of non-good Coxeter groups. The original excerpt also contains two TDLC-group questions omitted from the clean statement. Missing notation, finiteness assumptions, and the intended scope remain explicit. The supplied open status and dated literature assessment require independent review.
\subsection{Sources}
\begin{itemize}
\item Supplied local UnsolvedMath registry record, permanent ID \(30005708\), exact source problem code \texttt{catalogue entry}; title: ``Profinite Rigidity of Rational Euler Characteristics for Coxeter Groups.'' The record's clean statement and original excerpt have different scopes.
\item Supplied report citation: \emph{Mini-Workshop: Homological aspects of TDLC-groups} (2024), Oberwolfach Reports, pp.\ 3103--3138, DOI \texttt{10.4171/owr/2023/55}. Supplied link: \url{https://doi.org/10.4171/owr/2023/55}. The citation's 2024 year differs from the record's proposed year 2023 and its \texttt{owr-year:2023} tag; this discrepancy is unresolved.
\item Supplied attribution in the original excerpt: Questions 8--9, Yuri Santos Rego; Questions 10--11, Thomas Weigel; additional Questions 6--7, Roman Sauer. The structured \texttt{proposed\_by} field is null.
\item Supplied literature source, ``Oberwolfach report'': \url{https://ems.press/content/serial-article-files/48166}.
\item Supplied literature source, ``Hughes--Moller--Varghese'': \url{https://arxiv.org/abs/2505.08701}. No title or precise supporting result is supplied.
\item Supplied literature source, ``Kammeyer--Kionke--Wilkes'': \url{https://arxiv.org/abs/1901.01227}. No title or precise supporting result is supplied.
\item Literature assessment dated exactly \texttt{2026-08-22}: the supplied curation labels the problem open, states that finite-rank Coxeter groups are known to be virtually RFRS, and states that a Coxeter pair with isomorphic profinite completions and different rational Euler characteristics has neither been produced nor ruled out. This is supplied dated curation, not an independently verified literature assessment. The background explicitly requests independent verification of time-sensitive status and cited claims.
\item The record labels its statement \texttt{corrected\_verified} and reports comparison with the official report. This is supplied verification metadata; no source was visited or independently checked for this draft. The excerpt's attribution to Kielak under reference label [2] lacks a bibliographic entry in the supplied record.
\end{itemize}
\end{document}
